\documentclass[11pt,a4paper,landscape]{article}
\usepackage[T1]{fontenc}
\usepackage[utf8]{inputenc}
\usepackage[ngerman]{babel}
\usepackage{lmodern}
\usepackage{graphicx}
\usepackage{xcolor}
\usepackage{geometry}
\usepackage{array}
\usepackage{tabularx}
\usepackage{booktabs}
\usepackage{microtype}
\geometry{left=12mm,right=12mm,top=10mm,bottom=10mm}
\setlength{\parindent}{0pt}
\pagestyle{empty}
\definecolor{headblue}{RGB}{34,91,127}
\definecolor{lightblue}{RGB}{235,245,250}
\definecolor{goodgreen}{RGB}{62,130,70}
\definecolor{lightgreen}{RGB}{237,248,239}
\definecolor{lightgray}{RGB}{246,246,246}

\newcommand{\titleline}[1]{%
  {\color{headblue}\LARGE\bfseries #1}\par\vspace{1.5mm}%
  {\color{headblue}\rule{\textwidth}{1.2pt}}\par\vspace{3mm}%
}
\newcommand{\monkey}[2][3.55cm]{%
  \begin{minipage}[t]{4.0cm}\centering
    \includegraphics[width=#1,height=#1,keepaspectratio]{images/monkey_#2.png}\\[-1mm]
    {\large\bfseries Nr.~#2}
  \end{minipage}%
}
\newcommand{\monkeysmall}[2][3.15cm]{%
  \begin{minipage}[t]{3.55cm}\centering
    \includegraphics[width=#1,height=#1,keepaspectratio]{images/monkey_#2.png}\\[-1mm]
    {\normalsize\bfseries Nr.~#2}
  \end{minipage}%
}
\newcommand{\sourcefooter}{%
  \vfill\scriptsize\color{gray}%
  Affenbilder und Zuordnung der Testdaten gemäß der bereitgestellten AI-Unplugged-Präsentation
  „Klassifikation mit Entscheidungsbäumen: Das Gute-Äffchen-Böse-Äffchen-Spiel“ (Annabel Lindner, Stefan Seegerer).
}

\begin{document}

% --- Easy: student version ---
\titleline{Testdaten -- Einfache Variante}
\textbf{Schülerversion:} Die tatsächlichen Klassen sind noch nicht angegeben. Lass deinen fertigen
Entscheidungsbaum jedes Äffchen klassifizieren und notiere zunächst nur die Vorhersage. Vergleiche erst
anschließend mit der Auflösung auf der nächsten Seite.\par\vspace{5mm}

\begin{center}
\monkey{03}\hfill\monkey{05}\hfill\monkey{10}\hfill\monkey{11}\\[6mm]
\monkey{13}\hfill\monkey{16}\hfill\monkey{19}\hfill\monkey{20}
\end{center}
\sourcefooter
\newpage

% --- Easy: labels ---
\titleline{Testdaten -- Einfache Variante: Auflösung}
\begin{minipage}[t]{0.485\textwidth}
  \colorbox{lightgreen}{\parbox{0.96\linewidth}{\centering\Large\bfseries\color{goodgreen}Beißt}}
  \vspace{4mm}
  \begin{center}
    \monkey{03}\hfill\monkey{05}\\[5mm]
    \monkey{11}\hfill\monkey{19}
  \end{center}
\end{minipage}\hfill
\begin{minipage}[t]{0.485\textwidth}
  \colorbox{lightblue}{\parbox{0.96\linewidth}{\centering\Large\bfseries\color{headblue}Beißt nicht}}
  \vspace{4mm}
  \begin{center}
    \monkey{10}\hfill\monkey{13}\\[5mm]
    \monkey{16}\hfill\monkey{20}
  \end{center}
\end{minipage}
\sourcefooter
\newpage

% --- Advanced: student version ---
\titleline{Testdaten -- Fortgeschrittene Variante}
\textbf{Schülerversion:} Die tatsächlichen Klassen sind noch nicht angegeben. Klassifiziere die Äffchen
nacheinander mit deinem Entscheidungsbaum. Die Reihenfolge entspricht der bereitgestellten Präsentation.\par\vspace{3mm}

\begin{center}
\monkeysmall{03}\hfill\monkeysmall{06}\hfill\monkeysmall{08}\hfill\monkeysmall{11}\hfill\monkeysmall{13}\hfill\monkeysmall{18}\hfill\monkeysmall{20}\\[4mm]
\monkeysmall{26}\hfill\monkeysmall{27}\hfill\monkeysmall{29}\hfill\monkeysmall{31}\hfill\monkeysmall{34}\hfill\monkeysmall{21}
\end{center}
\sourcefooter
\newpage

% --- Advanced: labels ---
\titleline{Testdaten -- Fortgeschrittene Variante: Auflösung}
\begin{minipage}[t]{0.34\textwidth}
  \colorbox{lightgreen}{\parbox{0.95\linewidth}{\centering\Large\bfseries\color{goodgreen}Beißt}}
  \vspace{4mm}
  \begin{center}
    \monkeysmall{06}\hfill\monkeysmall{13}\\[5mm]
    \monkeysmall{18}\hfill\monkeysmall{34}
  \end{center}
\end{minipage}\hfill
\begin{minipage}[t]{0.63\textwidth}
  \colorbox{lightblue}{\parbox{0.97\linewidth}{\centering\Large\bfseries\color{headblue}Beißt nicht}}
  \vspace{4mm}
  \begin{center}
    \monkeysmall{03}\hfill\monkeysmall{08}\hfill\monkeysmall{11}\\[4mm]
    \monkeysmall{20}\hfill\monkeysmall{26}\hfill\monkeysmall{27}\\[4mm]
    \monkeysmall{29}\hfill\monkeysmall{31}\hfill\monkeysmall{21}
  \end{center}
\end{minipage}
\sourcefooter

\end{document}
